% ============================================================
% Come si usano, argomento 15: la gaussiana e la funzione Q
% ============================================================
\argomento{La variabile gaussiana e la funzione $Q$}{Formule a p.~\pageref*{f:gaussstd}}[Gaussiana e funzione Q]

\begin{scelta}{Come si calcola una probabilità gaussiana}
Ogni probabilità di una gaussiana $\mathcal{N}(\mu_X, \sigma_X^2)$ si riporta alla funzione $Q$ in tre passi:
(1) si \textbf{standardizza} la soglia, $x \mapsto \frac{x - \mu_X}{\sigma_X}$ (di quante deviazioni standard sta
sopra la media); (2) si scrive la probabilità con $Q$; (3) se l'argomento è negativo si usa $Q(-x) = 1 - Q(x)$ e si
legge la tavola \vv{tavolaq}.
\begin{center}
\renewcommand{\arraystretch}{1.45}
\begin{tabular}{@{}ll@{}}
  \toprule
  \textbf{Si chiede\ldots} & \textbf{si scrive\ldots} \\
  \midrule
  $P(X > x)$ & $Q\left(\frac{x - \mu_X}{\sigma_X}\right)$ \vv{codegauss} \\
  $P(X \leq x)$ & $1 - Q\left(\frac{x - \mu_X}{\sigma_X}\right)$ \\
  $P(a < X \leq b)$ & $Q\left(\frac{a - \mu_X}{\sigma_X}\right) - Q\left(\frac{b - \mu_X}{\sigma_X}\right)$ \\
  $P(\abs{X - \mu_X} \leq k\sigma_X)$, «entro $k$ sigma» & $1 - 2Q(k)$ \vv{tresigma} \\
  probabilità d'errore con livelli $\pm A$ e rumore gaussiano & $Q(A/\sigma)$ \vv{erroregauss} \\
  \bottomrule
\end{tabular}
\end{center}
\end{scelta}

\scheda{gaussstd}{
  \usoquando{È il riferimento: tutte le gaussiane si riportano alla standard, e le sue probabilità sono la funzione
    $Q$.}
  \usoesempio{$P(X_0 > 1) = Q(1) \simeq 0{,}159$; $P(-1 \leq X_0 \leq 1) = 1 - 2\,Q(1) \simeq 0{,}683$.}
}

\scheda{gaussintegrale}{
  \usoquando{Per verificare che la densità gaussiana abbia area 1, e negli integrali in cui compare $e^{-x^2/2}$ su
    tutta la retta.}
  \usocome{Si calcola il quadrato dell'integrale come integrale doppio, in coordinate polari (\vv{polari}).}
  \usoesempio{Con il cambio di variabile $u = x/\sigma$: $\int_{-\infty}^{\infty} e^{-\frac{x^2}{2\sigma^2}} \dd x =
    \sigma\sqrt{2\pi}$.}
}

\scheda{gaussgen}{
  \usoquando{Grandezze che sono la somma di tanti piccoli contributi indipendenti: rumore termico, errori di misura;
    oppure il testo dice «gaussiana» o «normale».}
  \usocome{Bastano media e varianza per conoscerla tutta. Per le probabilità si standardizza.}
  \usoattenzione{Il secondo parametro di $\mathcal{N}(\mu, \sigma^2)$ è la \emph{varianza}: in $\mathcal{N}(10, 4)$ la
    deviazione standard è 2, non 4.}
  \usoesempio{$X \sim \mathcal{N}(10, 4)$: $X_0 = \frac{X - 10}{2}$ è standard.}
}

\scheda{funzioneq}{
  \usoquando{Ogni volta che serve la probabilità che una gaussiana superi una soglia: la CDF gaussiana non ha una
    formula esplicita, e si usa la tavola di $Q$.}
  \usoesempio{$Q(2) \simeq 0{,}0228$: una gaussiana standard supera 2 circa 2 volte su cento.}
}

\scheda{codegauss}{
  \usoquando{$X$ è gaussiana con media e varianza qualsiasi e si chiede la probabilità di una soglia o di un
    intervallo.}
  \usocome{Si standardizzano le soglie e si usa $Q$; per gli argomenti negativi, la simmetria.}
  \usoesempio{$X \sim \mathcal{N}(10, 4)$: $P(X > 13) = Q(1{,}5) \simeq 0{,}0668$;
    $P(X \leq 7) = 1 - Q(-1{,}5) = Q(1{,}5) \simeq 0{,}0668$; $P(8 < X \leq 12) = Q(-1) - Q(1) = 1 - 2Q(1) \simeq
    0{,}683$.}
}

\scheda{qproprieta}{
  \usoquando{Per gli argomenti negativi (la tavola si ferma a $x \geq 0$) e per controllare i conti: $Q$ scende da 1 a
    0 e vale $\frac{1}{2}$ in 0.}
  \usoesempio{$Q(-0{,}5) = 1 - Q(0{,}5) \simeq 1 - 0{,}309 = 0{,}691$.}
}

\scheda{tavolaq}{
  \usoquando{Per leggere i valori di $Q$ negli esercizi.}
  \usocome{Si arrotonda l'argomento al decimo; se serve più precisione si interpola fra due valori vicini.}
  \usoesempio{$Q(1{,}3) \simeq 0{,}0968$; $Q(2{,}5) \simeq 0{,}00621$.}
}

\scheda{qmaggiorazione}{
  \usoquando{Per valori grandi dell'argomento, fuori dalla tavola, o per dimostrare che una probabilità d'errore va
    a zero molto in fretta.}
  \usoesempio{$Q(3) < \frac{e^{-4{,}5}}{3\sqrt{2\pi}} \simeq 0{,}00148$, contro il valore vero $0{,}00135$: già a $x = 3$
    l'approssimazione è buona.}
}

\scheda{tresigma}{
  \usoquando{Si chiede la probabilità di restare entro $k$ deviazioni standard dalla media (o di uscirne).}
  \usoattenzione{Vale per la gaussiana. Per una variabile qualsiasi si può dire solo quanto garantisce Chebyshev
    (\vv{chebyshev}): almeno $0{,}75$ entro $2\sigma$, invece di $0{,}954$.}
  \usoesempio{Entro $2\sigma$: $1 - 2Q(2) \simeq 1 - 0{,}0455 = 0{,}954$; fuori da $3\sigma$: $2Q(3) \simeq
    0{,}0027$, meno di 3 volte su mille.}
}

\scheda{erroregauss}{
  \usoquando{Un bit è trasmesso con due livelli $\pm A$ equiprobabili, il rumore è gaussiano a media nulla e il
    ricevitore decide con la soglia 0.}
  \usocome{Sapendo il bit, il segnale ricevuto è gaussiano centrato in $\pm A$; l'errore è la coda oltre la soglia, e
    per simmetria i due casi sono uguali.}
  \usoesempio{$A/\sigma = 3$: $P_e = Q(3) \simeq 0{,}00135$. Raddoppiando $A/\sigma$ da 2 a 4 gli errori scendono da
    $0{,}0228$ a $3{,}2 \cdot 10^{-5}$.}
}
